\phantomsection
\color{unimain}\section*{\centering{\large{ПЕРЕЧЕНЬ СОКРАЩЕНИЙ}}}
\addcontentsline{toc}{section}{Перечень сокращений}
{\color{white}\fontsize{0.1pt}{0.1pt}\selectfont<begABBRS>}
\color{black}
\begin{longtable}{>{\raggedright\arraybackslash}m{2cm}>{\raggedright\arraybackslash}m{0.5cm}>{\raggedright\arraybackslash}m{20cm}}
\endfirsthead\endhead\endfoot\endlastfoot
GOOSE & -- & Generic Object-Oriented Substation Event; \\
IRF & -- & Internal Relay Fault; \\
PPS & -- & Pulse Per Second; \\
PTP & -- & Precision Time Protocol (протокол точного времени); \\
SNTP & -- & Simple Network Time Protocol (простой сетевой протокол времени); \\
АПВ & -- & автоматическое повторное включение; \\
АСУ~ТП & -- & автоматизированная система управления технологическими процессами; \\
АЦП & -- & аналого-цифровой преобразователь; \\
БВКЗ & -- & блокировка при внешних коротких замыканиях; \\
БИ & -- & блок испытательный; \\
БУ & -- & блок управления; \\
ВН & -- & высшее напряжение; \\
ГЗ & -- & газовая защита; \\
ДЗТ & -- & дифференциальная защита трансформатора; \\
ДО & -- & дочерние общества; \\
ДТЗт & -- & дифференциальная токовая защита с торможением; \\
ДТО & -- & дифференциальная токовая отсечка; \\
ДТм & -- & датчик температуры масла; \\
ДТо & -- & датчик температуры обмотки; \\
ДУм & -- & датчик уровня масла; \\
ЗАПВ & -- & запрет автоматического повторного включения; \\
ЗДЗ & -- & защита от дуговых замыканий; \\
ЗИП & -- & запасные части, инструменты и принадлежности; \\
ЗП & -- & защита от перегрузки; \\
ЗПО & -- & защита от потери охлаждения; \\
ИО & -- & измерительный орган / избирательный орган; \\
ИЧМ & -- & интерфейс человек-машина; \\
КЗ & -- & короткое замыкание; \\
КИ & -- & контроль изоляции; \\
КСО & -- & камеры сборные одностороннего обслуживания; \\
КЦТ & -- & контроль цепей тока; \\
ЛО & -- & логика отключения; \\
МЭК & -- & международная электротехническая комиссия; \\
НЗ & -- & нормально закрытый; \\
НКУ & -- & низковольтное комплектное устройство; \\
НН & -- & низшее напряжение; \\
НО & -- & нормально открытый; \\
НП & -- & нулевая последовательность; \\
ОВ & -- & оперативный вывод / обходной выключатель; \\
ОЗУ & -- & оперативное запоминающее устройство; \\
ОК & -- & отсечной клапан; \\
ОП & -- & обратная последовательность; \\
ОСФ & -- & орган сравнения фаз; \\
ОТ & -- & оперативный ток; \\
ПДС & -- & преобразователь дискретных сигналов; \\
ПЗУ & -- & постоянное запоминающее устройство; \\
ПК & -- & предохранительный клапан; \\
ПО & -- & пусковой орган / программное обеспечение; \\
ПС & -- & предупредительная сигнализация; \\
РАС & -- & регистратор аварийных событий; \\
РД & -- & реле давления; \\
РЗ & -- & релейная защита; \\
РЗА & -- & релейная защита и автоматика; \\
РПН & -- & устройство регулирования напряжения под нагрузкой; \\
РТПО & -- & реле (орган) тока пуска охлаждения; \\
РУ & -- & распределительное устройство; \\
РЭ & -- & руководство по эксплуатации; \\
СН & -- & среднее напряжение; \\
СО & -- & система охлаждения; \\
СС & -- & сборка сигналов; \\
ТЗ & -- & технологические защиты; \\
ТК & -- & токовый контроль; \\
ТН & -- & трансформатор напряжения; \\
ТО & -- & токовая отсечка; \\
ТС & -- & технологическая сигнализация; \\
ТТ & -- & трансформатор тока; \\
ТУ & -- & телеуправление/телеускорение; \\
УРОВ & -- & устройство резервирования при отказе выключателя; \\
ФК & -- & функциональная клавиша; \\
ФСУ & -- & функциональная схема устройства; \\
ЦП & -- & центральный процессор; \\
ШЭТ & -- & шкаф электротехнический типовой; \\
ЭМО & -- & электромагнит отключения. \\
{\color{white}\fontsize{0.1pt}{0.1pt}\selectfont<endABBRS>}
\end{longtable}
